% ===================================================================
%  TABULKA č. 6 — Dům uvnitř, nábytek, stravování, osvětlení.
%  Traduction 2026 d'apres texto/io, controlee sur texto/fr, texto/ru
%  et texto/uk. Consigne : voir texto/cs/10-tabulka-01.tex.
%
%  LE PLUS LONG TABLEAU DE LA COLONNE, cinq pieces et deux cent
%  quarante objets, et donc le second banc d'essai apres le corps du
%  tableau 2. Il donne le meme resultat : tout ce qui se manie tous
%  les jours a son mot tcheque — « koště » le balai, « kohoutek » le
%  robinet, « pánev » la poele, « dřez » l'evier, « prostěradlo » le
%  drap, « hřeben » le peigne, « sirky » les allumettes, « nůžky »
%  les ciseaux, « polštář » l'oreiller.
%
%  ET LE MOT LE PLUS BANAL DE LA CUISINE SEPARE ICI QUATRE LANGUES.
%  Le galicien disait « mistos », l'espagnol « cerillas », le
%  portugais « fósforos » ; le tcheque dit « sirky », qui n'est aucun
%  des trois et qui vient du soufre. Quatre langues, quatre mots,
%  pour l'objet le plus commun qui soit.
%
%  L'ECART DECLARE DU BLOC c1-06-1, LE PREMIER DES TROIS. Le
%  fac-simile ido y ouvre sur « (6) », qui est le savon de l'alinea
%  2 ; c'est la femme de chambre « (16) », comme le porte le francais
%  et comme la planche le grave. La colonne le rend, avec les douze
%  qui la precedent, et renvoji.py le passe sans rien dire.
% ===================================================================

%%K t06-apar-1 apar
\VUsaut{6.0mm}
\VUtitre{15.0pt}{60}{TABULKA č. 6}
\VUsaut{1.4mm}
\VUtitre{11.4pt}{0}{Dům uvnitř, nábytek, stravování,}
\VUsaut{0.4mm}
\VUtitre{11.4pt}{0}{osvětlení.}
\VUsaut{2.2mm}

%%K t06-apar-2 apar
Dům je dokončen. Janův strýc se zabydlel ve svém novém bytě, jehož uspořádání známe. Podívejme
se nyní, jak si dům zařídil nábytkem. Prohlédneme si nejprve:

%%K t06-tit-1 sub
\VUpk{10.2pt}{40}{Ložnice.}
\VUsaut{3.0mm}

%%K t06-c1-01-1 p
1. --- Přicházíme nevhod, protože pokojská je zaměstnána úklidem \VUgras{pokoje.}

%%K t06-c1-02-1 p
2. --- Na \VUgras{toaletním stolku} \textsuperscript{(1)} jsou sem tam různé předměty, jimiž
se člověk myje, češe, zkrátka upravuje. Jsou to \VUgras{umyvadlo} \textsuperscript{(2)} a
\VUgras{džbán na vodu} \textsuperscript{(3)}, \VUgras{kartáč na vlasy} \textsuperscript{(4)},
\VUgras{hřeben} \textsuperscript{(5)},
\VUgras{mýdlo} \textsuperscript{(6)} a \VUgras{houba} \textsuperscript{(7)}, konečně
čistá \VUgras{lahvička} \textsuperscript{(8)} na tekutou zubní vodu.

%%K t06-c1-03-1 p
3. --- Na zemi, před toaletním stolkem, zde je \VUgras{kbelík} \textsuperscript{(9)} na
sbírání špinavé vody.

%%K t06-c1-04-1 p
4. --- V \VUgras{předsíni} \textsuperscript{(10)} vidíme několik \VUgras{věšáků}
\textsuperscript{(13)} a dva \VUgras{deštníky} \textsuperscript{(11)} ve
\VUgras{stojanu na deštníky} \textsuperscript{(12)}.

%%K t06-c1-05-1 p
5. --- Na druhé straně dveří je \VUgras{zrcadlová skříň} \textsuperscript{(14)}, v níž je
\VUgras{prádlo} \textsuperscript{(15)}.

%%K t06-c1-06-1 p
6. --- Využijme toho, že \VUgras{pokojská} \textsuperscript{(16)} stele \VUgras{postel}
\textsuperscript{(17)}, a prohlédněme si její různé části. \VUgras{Rám postele}
\textsuperscript{(20)} obsahuje dobrou \VUgras{pružinovou matraci} \textsuperscript{(21)}, na
kterou Justýna hned položí vlněnou \VUgras{matraci} \textsuperscript{(22)}. Potom vezme ze
židlí dvě \VUgras{prostěradla} \textsuperscript{(23)} a \VUgras{přikrývku}
\textsuperscript{(24)}. Pak položí do čela postele \VUgras{příčný polštář}
\textsuperscript{(25)} a \VUgras{podhlavník} \textsuperscript{(26)}, které jsou s
\VUgras{peřinou} \textsuperscript{(27)} na
\VUgras{přehozu} \textsuperscript{(32)}. Prachovkou oprašuje \VUgras{záclony}
\textsuperscript{(19)} a \VUgras{nebesa} \textsuperscript{(18)}, a část práce je hotova.

%%K t06-c1-07-1 p
7. --- Potom bude muset trochu uklidit \VUgras{noční stolek} \textsuperscript{(28)}, natáhnout
\VUgras{budík} \textsuperscript{(29)}, připravit \VUgras{noční lampičku}
\textsuperscript{(30)}, vyjmout \VUgras{svíčku} \textsuperscript{(31)}, aby vyčistila svícen.

%%K t06-tit-2 sub
\VUsaut{5.0mm}
\VUpk{10.2pt}{40}{Košík na šití a krbové náčiní.}
\VUsaut{3.0mm}

%%K t06-c1-08-1 p
8. --- \VUgras{Košík na šití} \textsuperscript{(34)} blízko \VUgras{stoličky}
\textsuperscript{(33)} také velmi potřebuje uklidit. Bude muset složit \VUgras{výšivky}
\textsuperscript{(35)}, zavřít do \VUgras{pracovního stolku} \textsuperscript{(38)}
\VUgras{náprstek,} \VUgras{nit} \textsuperscript{(36)}, \VUgras{jehly}
\textsuperscript{(37)} a \VUgras{nůžky} \textsuperscript{(39)} a zamknout
\VUgras{klíčem} \textsuperscript{(40)} víko stolku.

%%K t06-c1-09-1 p
9. --- Na \VUgras{jednonohém stolku} \textsuperscript{(41)} uprostřed Justýna obnoví vodu v
\VUgras{karafě} \textsuperscript{(42)}. Nasype \VUgras{cukr} do
\VUgras{cukřenky} \textsuperscript{(43)}, vyčistí \VUgras{kleštičky na cukr}
\textsuperscript{(44)} a odnese \VUgras{kartáč na šaty} \textsuperscript{(46)}.

%%K t06-c1-10-1 p
10. --- Pravá strana pokoje bude po ní žádat méně práce, protože \VUgras{lenoška}
\textsuperscript{(47)} a její \VUgras{polštář} \textsuperscript{(48)} jsou na svém správném
místě a, protože není zima, nic není přemístěno mezi náčiním \VUgras{krbu}
\textsuperscript{(49)}. \VUgras{Lopatka} \textsuperscript{(50)}, \VUgras{kleště}
\textsuperscript{(51)}, \VUgras{kozlíky} \textsuperscript{(55)} a \VUgras{zábrana na popel}
\textsuperscript{(56)} jsou čisté; \VUgras{krbová zástěna} \textsuperscript{(54)} nebyla
přemístěna a \VUgras{bedna na dříví} \textsuperscript{(52)} je dobře zásobena \VUgras{polínky}
\textsuperscript{(53)}.

%%K t06-noto-1 noto
\VUnotes{143.2mm}{%
(*) Babo = malé dítě; angl. \textit{baby}, fr. \textit{bébé}, it. \textit{bambino}.%
}

%%K t06-noto-2 noto
\VUnotes{143.2mm}{%
(*) Lambrequino = fr. \textit{lambrequin}, šp. \textit{lambrequines}, port.
\textit{lambrequins}, a dokonce něm. a angl. \textit{lambrequins}; tedy v němčině, angličtině,
francouzštině, portugalštině a španělštině.%
}

%%K t06-c1-11-1 p
11. --- Přesto bude muset oprášit \VUgras{kyvadlové hodiny} \textsuperscript{(57)},
\VUgras{svícny} \textsuperscript{(58)} a \VUgras{misky na drobnosti}
\textsuperscript{(59)}, konečně utřít \VUgras{zrcadlo} \textsuperscript{(60)}.

%%K t06-c1-12-1 p
12. --- Co se týče \VUgras{kolébky} \textsuperscript{(61)} děťátka (miminka (*)), ta je už
připravena.

%%K t06-tit-3 sub
\VUsaut{5.0mm}
\VUpk{10.2pt}{40}{Salon.}
\VUsaut{3.0mm}

%%K t06-c2-01-1 p
1. --- Nyní zde je salon. Je to velmi krásný \VUgras{pokoj.} Krb, který je v pravém rohu, je
ozdoben jemnou \VUgras{soškou} \textsuperscript{(1)}, dvěma krásnými
\VUgras{vázami} \textsuperscript{(3)} a zdoben sametovým \VUgras{lambrekýnem}
\textsuperscript{(2)} (*).

%%K t06-c2-02-1 p
2. --- Aby jiskry z ohniště nezapálily koberec, položili před ohniště měděnou
\VUgras{jiskrovou zástěnu} \textsuperscript{(4)}.

%%K t06-c2-03-1 p
3. --- Blízko je otevřený elegantní dámský \VUgras{psací stůl} \textsuperscript{(5)}. Na něm
píše \VUgras{paní domu} \textsuperscript{(23)} své dopisy.

%%K t06-c2-04-1 p
4. --- Okno je opatřeno \VUgras{záclonami} \textsuperscript{(6)} s \VUgras{úvazy}
\textsuperscript{(8)} zavěšenými na \VUgras{háčcích} \textsuperscript{(7)} a látkovými závěsy
\VUgras{nahoře nabranými} \textsuperscript{(9)}.

%%K t06-c2-05-1 p
5. --- Ve výklenku okna je v \VUgras{květináči} \textsuperscript{(11)} zelená
\VUgras{rostlina} \textsuperscript{(10)}, nádherná palma, jejíž listy sahají téměř až k
\VUgras{petrolejové lampě} \textsuperscript{(12)}.

%%K t06-c2-06-1 p
6. --- \VUgras{Stínítko} \textsuperscript{(13)} lampy upravila Marie, půvabná
\VUgras{dívka} \textsuperscript{(14)}, která je u klavíru \textsuperscript{(15)}. Sedí
velmi zpříma na \VUgras{stoličce} \textsuperscript{(16)} a studuje skladbu. Je velmi zručná a
pečlivá, jak dokazuje její \VUgras{skříň na noty} \textsuperscript{(17)}, která je dokonale
uspořádaná.

%%K t06-tit-4 sub
\VUsaut{5.0mm}
\VUpk{10.2pt}{40}{Nezbeda.}
\VUsaut{3.0mm}

%%K t06-c2-07-1 p
7. --- Její bratr Pavel hledá \VUgras{knihu} \textsuperscript{(19)} v
\VUgras{knihovně} \textsuperscript{(18)}. Je to \VUgras{mladík}
\textsuperscript{(20)} živý a nesnesitelný. Když se zachycuje knihovny, aby dosáhl na knihu,
která je příliš vysoko, riskuje, že shodí \VUgras{bustu} \textsuperscript{(21)}, která je
nahoře.

%%K t06-c2-08-1 p
8. --- Využívá k té hlouposti toho, že jeho matka je zaměstnána rozhovorem s přítelkyní,
\VUgras{dámou} \textsuperscript{(22)}, která přišla strávit několik dní u ní. Matka sedí na
\VUgras{pohovce} \textsuperscript{(24)} blízko \VUgras{křesla} \textsuperscript{(25)}, a její
přítelkyně je na \VUgras{sedadle} \textsuperscript{(27)}, z něhož je vidět jen
\VUgras{opěradlo} \textsuperscript{(26)}.

%%K t06-c2-09-1 p
9. --- Velký \VUgras{rám} \textsuperscript{(30)} visící na stěně obsahuje
\VUgras{podobiznu} \textsuperscript{(29)} jedné příbuzné. V \VUgras{albu}
\textsuperscript{(32)} položeném na stole jsou shromážděny \VUgras{fotografie}
\textsuperscript{(33)} ostatních členů rodiny. Děti si je právě prohlížely, protože
\VUgras{židle} \textsuperscript{(31)} jsou ještě seskupeny kolem stolu.

%%K t06-c2-10-1 p
10. --- \VUgras{Čínská váza} \textsuperscript{(34)} z porcelánu, která je blízko
\VUgras{nože na papír} \textsuperscript{(35)}, byla Marii darována k jejímu svátku, a
\VUgras{zástěnu} \textsuperscript{(36)}, která zdobí roh pokoje, přivezl z Číny její
strýc.

%%K t06-c2-11-1 p
11. --- Rozveselen překrásnými \VUgras{obrazy} \textsuperscript{(37)} zavěšenými na
\VUgras{příčce} \textsuperscript{(42)}, osvětlen \VUgras{stropním lustrem}
\textsuperscript{(38)}, jehož \VUgras{elektrické žárovky} \textsuperscript{(39)} večer vesele
září, ten salon vypadá velmi příjemně. Měkký \VUgras{koberec} \textsuperscript{(40)}
rozprostřený na parketách a tak snadno použitelný
\VUgras{elektrický zvonek} \textsuperscript{(41)} dokončují jeho pohodlí.

%%K t06-tit-5 sub
\VUsaut{3.6mm}
\VUpk{10.2pt}{40}{Jídelna.}
\VUsaut{3.0mm}

%%K t06-c3-01-1 p
1. --- Zazvonil zvonek k večeři. Celá rodina, kromě Marie, je shromážděna kolem stolu. Jídelna
je příjemně vytápěna výbornými \VUgras{kamny} \textsuperscript{(1)} z kameniny, s tenkou
\VUgras{rourou} \textsuperscript{(2)}.

%%K t06-c3-02-1 p
2. --- Ten pokoj neobsahuje mnoho nábytku. Podél zdi visí, nad \VUgras{policí}
\textsuperscript{(4)}, ozdobná \VUgras{kašnička} \textsuperscript{(3)}. Velmi blízko je
\VUgras{příborník} \textsuperscript{(5)}, na který se kladou studená jídla, zákusky a
různé stolní náčiní, jehož není hned třeba.

%%K t06-c3-03-1 p
3. --- Tak vidíme na horní desce \VUgras{pečené kuře} \textsuperscript{(7)},
\VUgras{rybu} \textsuperscript{(8)} a \VUgras{mísu} \textsuperscript{(10)} s
\VUgras{ovocem} \textsuperscript{(9)}. Dole, zde je \VUgras{sýr}
\textsuperscript{(11)}, \VUgras{chléb} \textsuperscript{(12)}, \VUgras{stojánek na lahvičky}
\textsuperscript{(13)}, který obsahuje vše potřebné k okořenění salátu: olej, ocet,
\VUgras{sůl} \textsuperscript{(14)} a \VUgras{pepř} \textsuperscript{(15)}. Konečně na spodní
desce je \VUgras{koláč} \textsuperscript{(17)} blízko \VUgras{hořčičnice}
\textsuperscript{(16)}.

%%K t06-c3-04-1 p
4. --- Hodiny v jídelně, kterým se říká \VUgras{nástěnné hodiny} \textsuperscript{(20)}, mají
velmi zvláštní tvar. Jsou vsazeny do dřevěného rámu, který obsahuje ještě \VUgras{tlakoměr}
\textsuperscript{(19)} a \VUgras{teploměr} \textsuperscript{(18)}.

%%K t06-tit-6 sub
\VUsaut{4.6mm}
\VUpk{10.2pt}{40}{Podávání večeře.}
\VUsaut{3.0mm}

%%K t06-c3-05-1 p
5. --- Na všech zdech pokoje jsou upevněna \VUgras{dřevěná ostění} \textsuperscript{(21)} z
voskovaného dubu. Látková \VUgras{dveřní záclona} \textsuperscript{(22)} dostatečně zavírá
dveře, které dávají vstup na verandu. Tam půjde rodina po večeři pít kávu. Už jsou
\VUgras{šálky} \textsuperscript{(24)} a
\VUgras{podšálky} \textsuperscript{(25)} připraveny na \VUgras{nádobníku}
\textsuperscript{(23)}.

%%K t06-c3-06-1 p
6. --- Nad stolem je ke stropu připevněna \VUgras{závěsná lampa} \textsuperscript{(26)}, jejíž
velmi jasné světlo večer osvětluje celou jídelnu.

%%K t06-c3-07-1 p
7. --- Nyní si prohlédněme samotný \VUgras{jídelní stůl} \textsuperscript{(27)}. Když rodina
vstoupila, příbor byl připraven. Na \VUgras{ubrusu} \textsuperscript{(28)} byl položen na
místě každého spolustolovníka \VUgras{talíř} \textsuperscript{(33)},
\VUgras{ubrousek} \textsuperscript{(29)}, \VUgras{lžíce} \textsuperscript{(34)},
\VUgras{vidlička} \textsuperscript{(35)}, \VUgras{nůž} \textsuperscript{(36)} a
\VUgras{sklenice} \textsuperscript{(32)}. Byly tam také \VUgras{láhve}
\textsuperscript{(30)} na víno a \VUgras{karafa} \textsuperscript{(31)} na vodu.

%%K t06-c3-08-1 p
8. --- \VUgras{Sluha} \textsuperscript{(39)} přinesl nejprve \VUgras{polévkovou mísu}
\textsuperscript{(37)}, v níž se ještě kouří trochu polévky. Nyní podává první
\VUgras{chod} \textsuperscript{(38)}, masitý chod. Potom předloží ty, které jsme už
jmenovali, konečně, po \VUgras{zákusku}, využije toho, že jeho páni vstoupí na verandu, aby
odnesl stolní náčiní a znovu uklidil jídelnu.

%%K t06-c3-08-2 p
\textit{Poznámka.} --- \textit{Běžná slova týkající se potravin jsou zde uvedena jen velmi
stručně. Seznam bude doplněn na dvou tabulkách «Hotel» (č. 13) a «Trh» (č. 15).}

%%K t06-tit-7 sub
\VUsaut{4.6mm}
\VUpk{10.2pt}{40}{Kuchyně.}
\VUsaut{3.0mm}

%%K t06-c4-01-1 p
1. --- V kuchyni, na kterou hledíme pootevřenými dveřmi, vidíme malebný nepořádek. Před
\VUgras{ohništěm} \textsuperscript{(1)}, v němž praská \VUgras{oheň}
\textsuperscript{(3)}, je postaven \VUgras{otáčecí rožeň} \textsuperscript{(5)}. Krásná
\VUgras{pečeně} \textsuperscript{(7)} je \VUgras{napíchnuta na rožni}
\textsuperscript{(6)} a dopéká se.

%%K t06-c4-02-1 p
2. --- \VUgras{Měch} \textsuperscript{(8)} a \VUgras{lopatka na uhlí} \textsuperscript{(9)},
kterých právě užili, zůstaly na zemi stejně jako
\VUgras{polínka} \textsuperscript{(10)}.

%%K t06-c4-03-1 p
3. --- Vršek krbu je uklizen; \VUgras{lucerna} \textsuperscript{(11)},
\VUgras{krabička na sirky} \textsuperscript{(13)}, v níž jsou \VUgras{sirky}
\textsuperscript{(12)}, \VUgras{svícen} \textsuperscript{(14)} a \VUgras{svícínek}
\textsuperscript{(15)} se \VUgras{svíčkou} \textsuperscript{(16)} jsou na svém místě. Ale
velký \VUgras{kbelík na uhlí} \textsuperscript{(18)}, plný \VUgras{černého uhlí}
\textsuperscript{(17)}, které \VUgras{kuchařka} \textsuperscript{(23)} právě naplnila ve
sklepě, zůstal uprostřed kuchyně.

%%K t06-c4-04-1 p
4. --- Vskutku, ubohá žena musí pospíchat, aby byl oběd hotov včas. Nyní je blízko
\VUgras{sporáku} \textsuperscript{(19)} a míchá \VUgras{omáčku}
\textsuperscript{(24)}, kterou nesmí příliš připálit.

%%K t06-c4-05-1 p
5. --- V \VUgras{troubě} \textsuperscript{(20)} se peče koláč, který připravila k svačině
dětí. Nad sporákem visí \VUgras{naběračka} \textsuperscript{(21)} a
\VUgras{pěnovačka} \textsuperscript{(22)}.

%%K t06-tit-8 sub
\VUsaut{4.0mm}
\VUpk{10.2pt}{40}{Nádobí.}
\VUsaut{3.0mm}

%%K t06-c4-06-1 p
6. --- \VUgras{Kuchyňské náčiní} \textsuperscript{(25)}, zavěšené vzadu v pokoji, je velmi
čisté. Krásné měděné \VUgras{kastroly} \textsuperscript{(26)}, \VUgras{konvice na mléko}
\textsuperscript{(27)}, \VUgras{cedník} \textsuperscript{(28)} a
\VUgras{struhadlo} \textsuperscript{(29)} byly pečlivě vyčištěny.

%%K t06-c4-07-1 p
7. --- \VUgras{Solnička} \textsuperscript{(30)} je položena na nádobníku blízko
\VUgras{čajové konvice} \textsuperscript{(31)} a \VUgras{velké láhve na olej}
\textsuperscript{(32)}. \VUgras{Zásuvka} \textsuperscript{(33)} pravděpodobně obsahuje příbory
a kuchyňské nože.

%%K t06-c4-08-1 p
8. --- \VUgras{Koště} \textsuperscript{(34)} a \VUgras{prachovka} \textsuperscript{(35)} jsou
v koutě. Kuchařka právě použila \VUgras{dřevěný špalek} \textsuperscript{(36)}, nechala jej
blízko okna.

%%K t06-c4-09-1 p
9. --- \VUgras{Ručník} \textsuperscript{(37)} visí na zdi, blízko \VUgras{nálevky}
\textsuperscript{(38)}. V té části kuchyně jsou uloženy \VUgras{rošt} \textsuperscript{(39)},
\VUgras{sekáček} \textsuperscript{(40)} a \VUgras{špalek na sekání} \textsuperscript{(41)}.
Potom, nad sekáčkem, na \VUgras{polici} \textsuperscript{(42)}, jsou \VUgras{hrnce}
\textsuperscript{(43)}, \VUgras{hrnec na vaření} \textsuperscript{(44)} se svou
\VUgras{poklicí} \textsuperscript{(45)} a
\VUgras{kotlík} \textsuperscript{(46)}. \VUgras{Pánev} \textsuperscript{(47)} visí
blízko.

%%K t06-c4-10-1 p
10. --- Jen měděný \VUgras{kohoutek} \textsuperscript{(48)} vrhá lesk do tohoto kouta.
\VUgras{Dřez} \textsuperscript{(49)}, na kterém je položen tlustý \VUgras{džbán}
\textsuperscript{(50)}, je pravděpodobně velmi pohodlný se svou malou skříňkou, v níž se
uchovává \VUgras{soudek na ocet} \textsuperscript{(51)} opatřený svým
\VUgras{kohoutkem} \textsuperscript{(52)}, \VUgras{bedna na odpadky}
\textsuperscript{(53)} a \VUgras{špinavé nádobí} \textsuperscript{(54)}.

%%K t06-tit-9 sub
\VUsaut{4.0mm}
\VUpk{10.2pt}{40}{Jiné předměty v kuchyni.}
\VUsaut{3.0mm}

%%K t06-c4-11-1 p
11. --- \VUgras{Škopek} \textsuperscript{(55)}, v němž se myje \VUgras{zelenina}
\textsuperscript{(58)}, a litinový \VUgras{kotel} \textsuperscript{(56)} položený na
\VUgras{dlaždicové podlaze} \textsuperscript{(57)} mají pravděpodobně také své místo v
té skříňce.

%%K t06-c4-12-1 p
12. --- Kuchařka jistě nemá nedostatek kamen, vždyť je tu ještě, na rohu stolu,
\VUgras{plynový vařič} \textsuperscript{(59)}, na němž se ohřívá voda v malém kastrolu.
\VUgras{Pára} \textsuperscript{(60)}, která z něho uniká, nám naznačuje, že
pravděpodobně vře. Ta voda se patrně použije na kávu, protože zde je na druhém konci stolu
\VUgras{konvice na kávu} \textsuperscript{(64)} a \VUgras{mlýnek na kávu}
\textsuperscript{(63)}.

%%K t06-c4-13-1 p
13. --- Vařič je spojen s \VUgras{plynovým hořákem} \textsuperscript{(62)} dlouhou
\VUgras{gumovou hadičkou} \textsuperscript{(61)}.

%%K t06-c4-14-1 p
14. --- \VUgras{Posluhovačka} \textsuperscript{(66)}, která přichází pomáhat kuchařce,
připravuje zeleninu na večeři. Až bude očištěna a umyta, položí ji do \VUgras{mísy}
\textsuperscript{(65)} a bude čekat na hodinu, kdy se bude vařit.

%%K t06-tit-10 sub
\VUsaut{4.0mm}
{\centering\fontsize{10.2pt}{10.2pt}\selectfont\textls[40]{\scshape Koupelna.} \textit{(Pokoj ke koupání.)}\par}
\VUsaut{3.0mm}

%%K t06-c5-01-1 p
1. --- Zbývá nám udělat jen jednu nediskrétnost, abychom poznali hlavní pokoje bytu: podívat
se na zařízení koupelny. To nebude trvat dlouho, protože není velká, a rychle zjistíme, že
\VUgras{vana} \textsuperscript{(1)} má dobrý \VUgras{ohřívač} \textsuperscript{(2)}, že
\VUgras{mýdlenka} \textsuperscript{(3)} je na dosah ruky koupajícího se a že
\VUgras{věšák na ručníky} \textsuperscript{(5)} jistě usnadní sušení \VUgras{ručníků}
\textsuperscript{(4)}.

%%K t06-c5-02-1 p
2. --- Ten pokoj má stěny pokryté \VUgras{kachlíky} \textsuperscript{(6)}, což umožnilo zřídit
\VUgras{sprchu} \textsuperscript{(7)} bez jakékoli obavy, že voda, která při jejím užívání
stříká zpět, poškodí zdi. Zinková \VUgras{mísa} \textsuperscript{(8)}, která slouží ke
koupelím nohou, dokončuje zařízení.
\VUsaut{6.0mm}
\VUfilet{22mm}
